% sec-01: Document 01, SB 1188, the California site authorization act.
% STAGE 3.  Two column widths retuned; finding (b) tightened by one clause;
% the seven-step authorization table moved to a page-breaking form, because a
% seven-row table with a four-line consequence cell is taller than the space the
% five subdivisions above it leave.
\docpart{SB 1188 Site Authorization}{SB 1188 - California PDAC LLM Oncology
Clinical Trial Site Authorization and Site Establishment Act of 2026}

\section{Legislative Findings and Declarations}

The Legislature finds and declares all of the following.

\begin{enumerate}
\item Pancreatic ductal adenocarcinoma remains among the most lethal solid
  tumors, and the resectable setting is served by fewer interventional trials
  than the metastatic setting. The one large readout available in 2026,
  RASolute 302, reported a median overall survival of 13.2 months against 6.6 in
  previously treated metastatic disease and is silent on the resectable setting
  \cite{rasolute302}.

\item Simulation work completed between 2025 and 2026 supports, without
  establishing, the hypothesis that a targeted agent combined with a staged
  robotic resection can be studied safely in the resectable setting. A ten-arm
  quantitative systems pharmacology run returned a median overall survival of
  12.8 months against 5.4 for chemotherapy, and a 1000-patient digital twin
  returned a best progression-free survival hazard ratio of 0.31
  \cite{qspmetpancre,daraxsims}.

\item That simulation work carries its own stated limitations, which the
  Legislature adopts rather than sets aside: the pharmacology run assumes no
  acquired resistance and ideal pharmacodynamics, and the digital twin uses no
  patient-specific pharmacokinetic, pharmacodynamic, or tumor growth parameters
  \cite{qspmetpancre,daraxsims}.

\item A credibility assessment aligned to ASME V\&V 40 and to ICH M15 scored
  81.9 across 55 verification notebook tests, which is a measured figure and not
  a regulatory finding \cite{daraxsims,asmevv40,ichm15}.

\item Three adapted regulatory frameworks covering good clinical practice, human
  subject protection, and investigational new drug applications have been
  published for physical artificial intelligence oncology trials and provide a
  workable governance spine for a site of this kind
  \cite{iche6r3adapt,cfr50adapt,cfr312adapt}.

\item A twenty-four-hour physical artificial intelligence oncology trial
  simulation demonstrated that a robot fleet can sustain a clinical schedule at
  99.7 percent uptime with all seven documented adverse events at Grade 1 or
  Grade 2 and all resolved within the hour they occurred \cite{newtrial2026}.

\item California has already legislated in adjacent territory. AB 3030 governs
  the use of generative artificial intelligence in patient communications,
  AB 489 addresses deceptive professional terms, SB 1120 governs utilization
  review, and AB 2013 requires training data transparency
  \cite{ab3030_2024,ab489_2025,sb1120_2024,ab2013_2024}. None of the four
  reaches the interventional clinical trial setting, and none addresses a model
  whose output is retained as a source record.

\item A large language model confined to advisory output does not practice
  medicine and does not replace a clinician. The distinction between advisory
  output and directive output is the distinction on which participant safety in
  this setting depends, and it is capable of being written as a testable rule
  \cite{cliniciantrust}.

\item An independent investigator, operating outside a university, can now
  assemble a complete regulatory package for a Phase 1 trial, and has done so
  \cite{phase1protocol,indapp}. State law should not be the reason such a site
  cannot open.

\item California has an opportunity to authorize the first such site in the
  nation. The authorization created by this chapter is conditional, revocable,
  time-limited, and reportable, because a first authorization that is none of
  those things is not an authorization but a concession.
\end{enumerate}

\section{Definitions}

For purposes of this chapter, the following definitions apply. A term defined
here carries the same meaning in documents 02, 03, 06, and 07 of this package.

\begin{enumerate}
\item ``LLM-directed clinical workflow'' means a sequence of clinical tasks in
  an interventional trial in which a large language model produces output that
  is presented to a licensed clinician before any action is taken on it.

\item ``Advisory-only output'' means output that cannot reach a participant,
  cannot alter a device setting, cannot change a dose, and cannot enter the
  clinical record except through an affirmative act by a named human, recorded
  at the time it is taken.

\item ``Robotic procedure suite'' means a room in which one or more robot
  instances operate within a defined envelope during a study procedure.

\item ``Robot instance'' means one physical robot, identified by serial number,
  registered under §4 of the regulations adopted pursuant to this chapter.

\item ``Site safety officer'' means the individual responsible for robot
  envelopes, interlocks, and the emergency stop chain, who holds the training
  required by the standard operating procedures of the site.

\item ``Qualified site'' means a facility that has met the building, premises,
  parking, staffing, equipment, and regulatory requirements established under
  this chapter and holds a current authorization from the department.

\item ``Sponsor-investigator'' has the meaning given in 21 CFR 312.3, and the
  department shall not authorize a site at which a person who controls the
  sponsor also serves as a clinical investigator \cite{cfr312,cfr54}.

\item ``Audit trail'' means a hash-chained, append-only record of every model
  prompt, every model completion, and every human disposition of a completion,
  retained under §5 of document 03 of this package.

\item ``Department'' means the California Department of Public Health.
\end{enumerate}

\section{Site Authorization}

\begin{enumerate}
\item An applicant shall file an application containing the contents specified
  in the regulations adopted under this chapter. The department shall review the
  application for completeness within 30 days of filing. If the department does
  not act within 30 days, the application is deemed complete.

\item The department shall complete substantive review within 90 days of the
  date an application is complete. If the department does not act within 90
  days, a conditional authorization is deemed granted for a period of 180 days,
  during which the department may complete its review.

\item A conditional authorization is valid for 24 months and may be renewed for
  successive 24-month terms.

\item No participant may be enrolled at a conditionally authorized site until
  all of the following have occurred: institutional review board approval of the
  protocol; a safe-to-proceed determination on the investigational new drug
  application under 21 CFR 312.42 \cite{cfr312}; and a departmental site visit.

\item Full authorization may be granted after the site has treated six
  participants and has passed a departmental survey conducted under §5 of the
  regulations adopted under this chapter.
\end{enumerate}

\begin{mvltable}
\begin{xltabular}{\textwidth}{>{\raggedright\arraybackslash}p{0.9cm} >{\raggedright\arraybackslash}p{4.1cm} >{\raggedright\arraybackslash}p{2.2cm} >{\raggedright\arraybackslash}X}
\headrow \thead{Step} & \thead{Action} & \thead{Statutory clock} & \thead{Consequence if the department does not act}\\
\midrule
\endfirsthead
\headrow \thead{Step} & \thead{Action} & \thead{Statutory clock} & \thead{Consequence if the department does not act}\\
\midrule
\endhead
1 & Application filed & Day 0 & None; the clock starts \\
2 & Completeness review & 30 days & Application deemed complete \\
3 & Substantive review & 90 days & Conditional authorization deemed granted for 180 days \\
4 & Conditional authorization issued & 24 months & Expires; no enrollment may continue \\
5 & First participant hold released & On three conditions & Hold remains; no enrollment \\
6 & Departmental survey & After 6 participants & Conditional authorization continues to its term \\
7 & Full authorization & 24-month renewals & Site reverts to conditional status \\
\end{xltabular}
\end{mvltable}
\tabcapl{tab:sb1188auth}{The authorization sequence, with the consequence of
departmental silence stated for every step. A statute that is silent on
inaction has no schedule, only an aspiration.}

\section{Operating Conditions}

Each condition below is implemented by a document in this package, and the
department may verify compliance against that document.

\begin{mvtable}
\begin{tabularx}{\textwidth}{>{\raggedright\arraybackslash}p{4.6cm} >{\raggedright\arraybackslash}p{2.6cm} >{\raggedright\arraybackslash}X}
\headrow \thead{Condition} & \thead{Implemented by} & \thead{Evidence the department may request}\\
\midrule
Minimum staffing and delegation & Document 14 & The roster, the delegation of authority log, the training records \\
Premises zones and robot envelopes & Document 09 & The zone plan, the envelope survey, the interlock test record \\
Model governance and disclosure & Document 03 & The model version register and the audit trail extract \\
Safety reporting & Document 13 & The safety report log and the transmission receipts \\
Emergency response & Document 11 & The plan, the quarterly drill records, the after-action reports \\
Record retention & Document 03 & The retention schedule and the archive index \\
\end{tabularx}
\end{mvtable}
\tabcapl{tab:sb1188cond}{The six operating conditions and the artifact that
discharges each. No condition in this chapter is satisfied by a policy
statement; each is satisfied by a record.}

\section{Inspection, Enforcement, and Suspension}

\begin{enumerate}
\item The department may enter and inspect a qualified site during hours of
  operation without prior notice.

\item Deficiencies are classified as follows, and the same three classes are
  used in documents 03, 06, and 12 of this package so that one event is not
  graded on three scales.
  \begin{enumerate}
  \item Class A: a deficiency that presents immediate jeopardy to a participant.
    The department shall suspend the affected activity immediately, and the
    civil penalty is not less than \$10,000 and not more than \$25,000 per
    occurrence.
  \item Class B: a deficiency presenting a substantial probability of harm. The
    site shall file a correction plan within 10 days, and the civil penalty is
    not less than \$2,500 and not more than \$10,000.
  \item Class C: an administrative or record deficiency with no probability of
    harm. The site shall file a correction plan within 30 days, and the civil
    penalty is not more than \$2,500.
  \end{enumerate}

\item A site may appeal a classification to the department within 15 days, and
  the department shall decide the appeal within 45 days. A Class A suspension is
  not stayed by an appeal.
\end{enumerate}

\section{Reporting and Sunset}

\begin{enumerate}
\item The department shall report annually to the Legislature. The report shall
  state the number of authorized sites, the number of participants enrolled, the
  number and class of deficiencies cited, every model fault reported under §5 of
  document 11 of this package, and the department's recommendation on
  continuation.

\item This chapter shall remain in effect only until January 1, 2032, and as of
  that date is repealed, unless a later enacted statute deletes or extends that
  date. The annual report required by subdivision (a) is the record on which any
  extension is to be decided.
\end{enumerate}
